\pdfoutput=1

\documentclass[11pt]{article}

\usepackage[]{acl}
\usepackage{times}
\usepackage{latexsym}
\usepackage{hyperref}
\usepackage[T1]{fontenc}
\usepackage[utf8]{inputenc}
\usepackage{microtype}
\usepackage{booktabs}
\usepackage{graphicx}

\newcommand{\todo}[1]{\textcolor{orange}{[TODO: #1]}}

\title{Title of the Project\thanks{\hspace{0.5em}Some content is from Stanford CS224n. The \LaTeX\ template is adopted from ACL style and formatting.}}
\author{First Author \\
  UC Santa Cruz \\
  \texttt{email@ucsc.edu} \\\And
  Second Author \\
  UC Santa Cruz \\
  \texttt{email@ucsc.edu} \\\AND
  Third Author \\
  UC Santa Cruz \\
  \texttt{email@ucsc.edu} \\\And
  Fourth Author \\
  UC Santa Cruz \\
  \texttt{email@ucsc.edu} \\}

\begin{document}
\maketitle

\begin{abstract}
Your abstract should motivate the topic or problem and describe your goals. You should also highlight your main results and insights. Given that your project is still ongoing, it is okay if your results are still in progress. 
\end{abstract}

\section{Approach}
This section describes how you approach the problem. It should be clear and logical.
\begin{itemize}
    \item Please be specific when describing your main methods. For example, you may want to include necessary equations and figures (but it is OK to postpone developing figures until the final report).
    \item Describe your comparison methods. Depending on space limits and whether your baselines are standard, you may write this down in detail or refer to other papers for more information.
    \item Make it clear if any part of your approach is original. Provide references for models and techniques that are not yours.
    \item If you are using any code you did not write, make it clear and provide a reference or link. When describing something you implemented yourself, also make it clear.
\end{itemize}

\section{Experiments}
This section contains the following.
\begin{itemize}
    \item \textbf{Data:} Provide references for the dataset(s) you are using. Ensure the task associated with the dataset is clearly described.
    \item \textbf{Evaluation method:} Describe the evaluation metric(s) you adopted, as well as any other details required to understand your evaluation.
    \item \textbf{Experimental details:} Please describe how you ran your experiments (e.g., model configurations, sampling parameters, etc.).
    \item \textbf{Results:} Report your quantitative results so far. To compare several results or baseline methods, use a table or plot.
\end{itemize}
    
\section{Future Work}
Describe what you plan to do for the rest of your project and why.

\nocite{yao2022react}
\bibliographystyle{acl_natbib}
\bibliography{references}

\end{document}